%
% Section 9.6_RatioRootTest
%

%~~~~~~~~~~~~~~~~~~~~~~   []   ~~~~~~~~~~~~~~~~~~~~~~~~~~~
\begin{frame}
\frametitle{\Large {\color[rgb]{0.02,0.20,1.0} Section 9.6 : Ratio \& Root Test}}
\vspace{0.8cm}
\scriptsize
\begin{block}{Homework Problems:}

\vspace{0.2cm}
\centering
13-29(odd), 33, 37-43(odd), 47
	
\end{block}
\vspace{0.5cm}

\begin{block}{Section Example Problems to Read:}

\vspace{0.2cm}
\centering
1-5
	
\end{block}
\vspace{0.2cm}

\begin{block}{This section will enable you to ...}
%-------LIST---------
\begin{itemize}
	\item[...] apply the Ratio Test to determine the convergence or divergence of a series.
	\item[...] apply the Root Test to determine the convergence or divergence of a series.
\end{itemize}
%-------------------
\end{block}
\vspace{3.0cm}
\end{frame}

%~~~~~~~~~~~~~~~~~~~~~~~~~~~~~~~~~~~~~~~~~~~~~~~~~~~~~~~~~
\begin{frame}
\frametitle{I. Ratio Test}
\scriptsize
\begin{columns}
\begin{column}{12.0cm}
	\begin{columns}\fcolorbox{blue}{white}{\begin{column}{9.4cm}
	% Type the contents of the box here
	\underline{Ratio Test}:
	\quad Consider $\sum \limits_{n=1}^{\infty} b_n$ where $b_n$ is never zero. \\
%	\vspace{0.2cm}
\hspace{2.0cm} Let \qquad $\displaystyle{\lim \limits_{n \to \infty}\,\left|\frac{b_{n+1}}{b_{n}}\right|} = L_1$
	\begin{itemize}
		\item if \quad $L_1 < 1$ \quad then \quad $\sum \limits_{n=1}^{\infty} b_n$ \quad \onslide<2->{\textcolor{blue}{converges absolutely.}} \\
		\item if \quad $L_1 > 1$ \quad then \quad $\sum \limits_{n=1}^{\infty} b_n$ \quad \onslide<3->{\textcolor{blue}{diverges.}} \\
		\vspace{0.1cm}
		\item if \quad $L_1 = 1$ \quad then the ratio test is \onslide<4->{\textcolor{blue}{inconclusive.}} \\
	\end{itemize}
	%-------------------
	\end{column}}\begin{column}{0.0cm}\end{column}\end{columns}
	\vspace{0.1cm}
	Determine the convergence or divergence\\
	\vspace{0.1cm}
	$\displaystyle{\sum \limits_{n=1}^{\infty} \frac{2}{n!} \onslide<5->{\textcolor{blue}{\,= \,2\,\sum \limits_{n=1}^{\infty} \frac{1}{n!}}}}$
	\quad \onslide<5->{\textcolor{blue}{let \quad $\displaystyle{b_n = \frac{1}{n!}}$}}
	%--------EQN----------
	\begin{align*}
		\text{compute} \quad \onslide<6->{\textcolor{blue}{\lim \limits_{n \to \infty}\,\left|\frac{b_{n+1}}{b_{n}}\right| = \lim \limits_{n \to \infty}\,\left|b_{n+1}\cdot\frac{1}{b_{n}}\right|}} \onslide<7->{\textcolor{blue}{\, = \lim \limits_{n \to \infty}\,\left|\frac{1}{(n+1)!}\cdot \frac{n!}{1}\right|}} & \onslide<8->{\textcolor{blue}{\, = \lim \limits_{n \to \infty}\,\left|\frac{1}{(n+1)\only<1-8>{n!}\only<9->{\cancel{n!}}}\cdot\frac{\only<1-8>{n!}\only<9->{\cancel{n!}}}{1}\right|}}\\
		&\onslide<10->{\textcolor{blue}{\, = \lim \limits_{n \to \infty}\,\frac{1}{n+1}}} \onslide<11->{\textcolor{blue}{\, = 0 < 1}}
	\end{align*}
	%----------------------
	Therefore, by the ratio test \onslide<12->{\textcolor{blue}{\quad \fbox{$\displaystyle{\sum \limits_{n=1}^{\infty} \frac{2}{n!}}$ \quad converges absolutely.}}}
\end{column}
\begin{column}{0.0cm}\end{column}
\end{columns}
\end{frame}

%~~~~~~~~~~~~~~~~~~~~~~~~~~~~~~~~~~~~~~~~~~~~~~~~~~~~~~~~~
\begin{frame}
\frametitle{II. Root Test}
\scriptsize
\begin{columns}
\begin{column}{12.0cm}
	\begin{columns}\fcolorbox{blue}{white}{\begin{column}{9.4cm}
	% Type the contents of the box here
	\underline{Root Test}:
	\quad Consider $\sum \limits_{n=1}^{\infty} b_n$. \\
\hspace{2.0cm} Let \qquad $\displaystyle{\lim \limits_{n \to \infty}\,\sqrt[n]{|b_{n}|}} = L_2$
	%-------LIST---------
	\begin{itemize}
		\item[1.] if \quad $L_2 < 1$ \quad then \quad $\sum \limits_{n=1}^{\infty} b_n$ \quad \onslide<2->{\textcolor{blue}{converges absolutely.}} \\
		\item[2.] if \quad $L_2 > 1$ \quad then \quad $\sum \limits_{n=1}^{\infty} b_n$ \quad \onslide<3->{\textcolor{blue}{diverges.}} \\
		\vspace{0.1cm}
		\item[3.] if \quad $L_2 = 1$ \quad then the ratio test is \onslide<4->{\textcolor{blue}{inconclusive.}} \\
	\end{itemize}
	%-------------------
	\end{column}}\begin{column}{0.0cm}\end{column}\end{columns}
	\vspace{0.2cm}
	Determine the convergence or divergence\\
	\vspace{0.1cm}
	$\displaystyle{\sum \limits_{n=1}^{\infty} \frac{2^n}{n^n}}$
	\quad \onslide<5->{\textcolor{blue}{let \quad $\displaystyle{b_n = \frac{2^n}{n^n}}$}}
	%--------EQN----------
	\begin{align*}
		\text{compute} \quad \onslide<6->{\textcolor{blue}{\lim \limits_{n \to \infty}\,\left|\sqrt[n]{b_{n}}\right| = \lim \limits_{n \to \infty}\,\left|(b_n)^{1/n}\right|}} \onslide<7->{\textcolor{blue}{\, = \lim \limits_{n \to \infty}\,\left|\left(\frac{2^n}{n^n}\right)^{1/n}\right|}} & \onslide<8->{\textcolor{blue}{\, = \lim \limits_{n \to \infty}\,\left|\frac{\left(2^n\right)^{1/n}}{\left(n^n\right)^{1/n}}\right|}}\\
		&\onslide<9->{\textcolor{blue}{\, = \lim \limits_{n \to \infty}\,\frac{2}{n}}} \onslide<10->{\textcolor{blue}{\, = 0 < 1}}
	\end{align*}
	%----------------------
	Therefore, by the root test \onslide<11->{\textcolor{blue}{\quad \fbox{$\displaystyle{\sum \limits_{n=1}^{\infty} \frac{2^n}{n^n}}$ \quad converges absolutely.}}}
\end{column}
\begin{column}{0.0cm}\end{column}
\end{columns}
\end{frame}

%~~~~~~~~~~~~~~~~~~~~~~~~~~~~~~~~~~~~~~~~~~~~~~~~~~~~~~~~~
\begin{frame}
\frametitle{}
\scriptsize

\vspace{0.2cm}

\begin{columns}
\begin{column}{11.0cm}
Determine the convergence or divergence\\

$\displaystyle{\sum \limits_{n=1}^{\infty} \frac{(-1)^n n!}{2^n}}$ \qquad \onslide<2->{\textcolor{blue}{(Note: The factorial indicates you should try the ratio test)}}

\vspace{0.2cm}
\onslide<3->{\textcolor{blue}{Let \quad $\displaystyle{b_n = \frac{(-1)^n n!}{2^n}}$\qquad and to use the ratio test we need to compute $\displaystyle{\quad L = \lim \limits_{n \to \infty}\, \left|\frac{b_{n+1}}{b_n}\right|}$}}

\vspace{0.2cm}

%--------EQN----------
\begin{align*}
	\onslide<4->{\textcolor{blue}{L = \lim \limits_{n \to \infty}\, \left|\frac{b_{n+1}}{b_n}\right|}} &\, \onslide<5->{\textcolor{blue}{= \lim \limits_{n \to \infty}\, \left|\frac{(-1)^{n+1} (n+1)!}{2^{n+1}}\cdot \frac{2^n}{(-1)^n n!}\right|}}\\
	&\,\onslide<6->{\textcolor{blue}{= \lim \limits_{n \to \infty}\, \left|\frac{(-1)^{n+1}}{(-1)^{n}}\right| \cdot \left|\frac{(n+1)\only<7->{\cancel}{n!}}{\only<7->{\cancel}{2^n} \cdot 2^{1}}\cdot \frac{\only<7->{\cancel}{2^n}}{\only<7->{\cancel}{n!}}\right|}} \\
	&\,\onslide<8->{\textcolor{blue}{= \lim \limits_{n \to \infty}\, 1 \cdot \left|\frac{n+1}{2}\right|}} \\
	&\,\onslide<9->{\textcolor{blue}{= \lim \limits_{n \to \infty}\, \frac{n+1}{2} = \infty}}
\end{align*}
%----------------------

\vspace{0.2cm}

\onslide<10->{\textcolor{blue}{So $L > 1$ and by the ratio test \quad $\displaystyle{\sum \limits_{n=1}^{\infty} \frac{(-1)^n n!}{2^n}}$ \quad diverges.}}

\vspace{10.5cm}

\end{column}
\begin{column}{0.5cm}
	
\end{column}
\end{columns}
\end{frame}

%~~~~~~~~~~~~~~~~~~~~~~~~~~~~~~~~~~~~~~~~~~~~~~~~~~~~~~~~~
\begin{frame}
\frametitle{}
\scriptsize

\vspace{0.2cm}

\begin{columns}
\begin{column}{11.0cm}

$\displaystyle{\sum \limits_{n=1}^{\infty} \frac{n^{2n}}{(n+1)^n}}$ \qquad \onslide<2->{\textcolor{blue}{(Note: The powers of $n$ indicate you should try the root test)}}

\vspace{0.2cm}

\onslide<3->{\textcolor{blue}{Let \quad $\displaystyle{b_n = \frac{n^{2n}}{(n+1)^n}}$\qquad and to use the root test we compute $\displaystyle{\quad L = \lim \limits_{n \to \infty}\, \sqrt[n]{|b_n|}}$}}

\vspace{0.2cm}

%--------EQN----------
\begin{align*}
	\onslide<4->{\textcolor{blue}{L = \lim \limits_{n \to \infty}\, \sqrt[n]{|b_n|}}} \onslide<5->{\textcolor{blue}{= \lim \limits_{n \to \infty}\, |b_n|^{1/n}}}  &\, \onslide<6->{\textcolor{blue}{\,= \lim \limits_{n \to \infty}\, \left|\frac{n^{2n}}{(n+1)^n}\right|^{1/n}}} \\
	&\,\onslide<8->{\textcolor{blue}{= \lim \limits_{n \to \infty}\, \left|\frac{n^{2n/n}}{(n+1)^{n/n}}\right|}} \\
	&\,\onslide<9->{\textcolor{blue}{= \lim \limits_{n \to \infty}\, \frac{n^{2}}{(n+1)^{1}} \quad ^{\left(\frac{\infty}{\infty}\right)} }} \\
	&\,\onslide<10->{\textcolor{blue}{\overset{LH}{=} \lim \limits_{n \to \infty}\, \frac{2n}{1} = \infty}}
\end{align*}
%----------------------

\vspace{0.2cm}

\onslide<11->{\textcolor{blue}{So $L > 1$ and by the root test \quad $\displaystyle{\sum \limits_{n=1}^{\infty} \frac{n^{2n}}{(n+1)^n}}$ \quad diverges.}}

\vspace{10.5cm}

\end{column}
\begin{column}{0.5cm}
	
\end{column}
\end{columns}
\end{frame}

%~~~~~~~~~~~~~~~~~~~~~~~~~~~~~~~~~~~~~~~~~~~~~~~~~~~~~~~~~
\begin{frame}
\frametitle{}
\scriptsize

\vspace{0.2cm}

\begin{columns}
\begin{column}{11.0cm}

$\displaystyle{\sum \limits_{n=1}^{\infty} \frac{n\cdot2^n}{n!}}$ \qquad \onslide<2->{\textcolor{blue}{(Note: The factorial indicates you should try the ratio test)}}


\vspace{0.2cm}
\onslide<3->{\textcolor{blue}{Let \quad $\displaystyle{b_n = \frac{n\cdot2^n}{n!}}$\qquad and to use the ratio test we need to compute $\displaystyle{\quad L = \lim \limits_{n \to \infty}\, \left|\frac{b_{n+1}}{b_n}\right|}$}}

\vspace{0.2cm}

%--------EQN----------
\begin{align*}
	\onslide<4->{\textcolor{blue}{L = \lim \limits_{n \to \infty}\, \left|\frac{b_{n+1}}{b_n}\right|}} &\, \onslide<5->{\textcolor{blue}{= \lim \limits_{n \to \infty}\, \left|\frac{(n+1)\cdot2^{n+1}}{(n+1)!}\cdot \frac{n!}{n\cdot2^n}\right|}}\\
	&\,\onslide<6->{\textcolor{blue}{= \lim \limits_{n \to \infty}\, \left|\frac{\only<7->{\cancel}{(n+1)}\cdot\only<7->{\cancel}{2^n}\cdot 2^{1}}{\only<7->{\cancel}{(n+1)}\cdot\only<7->{\cancel}{n!}}\cdot \frac{\only<7->{\cancel}{n!}}{n\cdot\only<7->{\cancel}{2^n}}\right|}} \\
	&\,\onslide<8->{\textcolor{blue}{= \lim \limits_{n \to \infty}\, \left|\frac{2}{n}\right|}} \\
	&\,\onslide<9->{\textcolor{blue}{= \lim \limits_{n \to \infty}\, \frac{2}{n} = 0}}
\end{align*}
%----------------------

\vspace{0.2cm}

\onslide<10->{\textcolor{blue}{So $L < 1$ and by the ratio test \quad $\displaystyle{\sum \limits_{n=1}^{\infty} \frac{n\cdot2^n}{n!}}$ \quad converges.}}

\vspace{10.5cm}

\end{column}
\begin{column}{0.5cm}
	
\end{column}
\end{columns}
\end{frame}
